\documentclass[11pt]{article}
\usepackage[T1]{fontenc}
\usepackage{lmodern}
\usepackage[margin=1in]{geometry}
\usepackage{amsmath,amssymb,amsthm}
\usepackage{xcolor}
\usepackage[colorlinks=true,linkcolor=blue!45!black,urlcolor=blue!45!black,citecolor=blue!45!black]{hyperref}
\usepackage{microtype}
\setlength{\emergencystretch}{1em}
\newcommand{\Q}{\mathbb Q}
\newcommand{\Z}{\mathbb Z}
\newcommand{\F}{\mathbb F}
\newtheorem{theorem}{Theorem}
\newtheorem{lemma}[theorem]{Lemma}
\title{An odd-parameter factor exclusion\\for the second cuboid quintic}
\author{GPT-6 Astra}
\date{15 October 2026}
\begin{document}
\maketitle
\begin{abstract}
We give a Lean-checked proof that the normalized second cuboid quintic
has no rational divisor of degree two when its parameter is the square of
the ratio of two positive unequal odd integers. The proof transports factor
degrees through two finite-dimensional quotient-algebra isomorphisms, then
uses an elementary obstruction modulo four. It includes reducible and
repeated quadratic divisors and does not require coprimality of the parameters.
This is a research note on one parameter restriction. The opposite-parity
case and the full second cuboid conjecture remain unresolved here.
The originality of the restriction has not been established.
\end{abstract}

\section{Statement and scope}
For $s\in\Q$, define
\begin{align}
P_s(x)={}&x^5+(2+s)(3-2s)x^4 \notag\\
&+(1+10s+4s^2-14s^3+s^4)x^3 \notag\\
&-s(1-14s+4s^2+10s^3+s^4)x^2 \notag\\
&-s^3(1+2s)(-2+3s)x-s^5. \label{eq:P}
\end{align}
This is the normalization used in the recent quadratic-factor study
\cite{asiryan-rathbun}. For Sharipov's even degree-ten polynomial,
\[
Q_{p,q}(t)=q^{20}P_{(p/q)^2}(t^2/q^4).
\]
Our formal statement concerns $P_s$ itself; the displayed relation explains
its role in the factorization problem.

\begin{theorem}\label{thm:main}
Let $p,q\in\Z$ be positive, odd, and unequal, and put $s=(p/q)^2$.
There is no polynomial $D\in\Q[x]$ with $\deg D=2$ and $D\mid P_s$.
\end{theorem}

Coprimality is unnecessary. No irreducibility or squarefreeness assumption
is imposed on $D$. The theorem proves a restriction within the quadratic-factor
problem; it does not prove irreducibility of $P_s$ or $Q_{p,q}$.
For coprime positive parameters, the unhandled parity class has one parameter
even and the other odd. The condition $p\ne q$ is essential to the changes
of variable: $P_1(x)=(x-1)(x+1)^4$.

\section{Two changes of variable}
Throughout this section assume $s>0$ and $s\ne1$, and set
\[
h=\frac{s-1}{s+1},\qquad k=h^2.
\]
Then $h\ne0$ and $0<k<1$. Introduce the monic quintics
\begin{align*}
T_k(w)&=w(w^2-1)^2+4(1-k)(w^2-k),\\
g_k(v)&=v^5+v^4-2v^3-2v^2+2-2k.
\end{align*}

\begin{lemma}\label{lem:algebras}
There are $\Q$-algebra isomorphisms
\[
\Q[x]/(P_s)\ \cong\ \Q[w]/(T_k)\ \cong\ \Q[v]/(g_k).
\]
\end{lemma}
\begin{proof}
All three algebras have dimension five. In $A=\Q[x]/(P_s)$, the element
$x+s$ is a unit because
\[
P_s(-s)=-2s^3(s-1)^4\ne0.
\]
Define $w=h(x-s)/(x+s)$, with division interpreted as multiplication by
the inverse of a unit. The polynomial identity
\begin{equation}\label{eq:mobius}
(h-w)^5P_s\!\left(s\frac{h+w}{h-w}\right)
=2s^3(s-1)^4T_{h^2}(w)
\end{equation}
holds wherever the displayed denominators are defined. Clearing denominators
also gives an identity in the quotient algebra, so $T_k(w)=0$ in $A$.
Here $h-w=2hs/(x+s)$ is a unit. The inverse expression
\[
x=s\frac{h+w}{h-w}
\]
shows that $w$ generates $A$. To justify the use of the inverse within the
subalgebra generated by $w$, use the following elementary fact: in a
finite-dimensional commutative algebra over a field, a subalgebra containing
an ambient unit also contains its inverse. Multiplication by that unit is
injective on the subalgebra, hence surjective. Thus the homomorphism from
$\Q[w]/(T_k)$ to $A$ is surjective. Equal dimensions make it an isomorphism.

For the second isomorphism, work in $B=\Q[v]/(g_k)$ and set $w=1-v^2$.
Direct polynomial identities give
\begin{align}
T_k(1-v^2)&=g_k(v)g_k(-v), \label{eq:composition}\\
g_k(v)&=v^2(v^2-2)(v+1)+2(1-k), \label{eq:unit}\\
2k-1-(1-v^2)^2-v\big((1-v^2)^2-1\big)&=-g_k(v). \label{eq:inverse}
\end{align}
Equation~\eqref{eq:composition} gives $T_k(w)=0$ in $B$.
By~\eqref{eq:unit}, the element
$w^2-1=v^2(v^2-2)$ is a unit, since its product with $v+1$ is
$-2(1-k)\ne0$. Equation~\eqref{eq:inverse} therefore yields
\[
v=\frac{2k-1-w^2}{w^2-1}.
\]
The inverse-in-subalgebra argument again proves that $w$ generates $B$.
The resulting surjection $\Q[w]/(T_k)\longrightarrow B$ is an isomorphism.
\end{proof}

\begin{lemma}\label{lem:degree}
If $P_s$ has a degree-two divisor over $\Q$, then $g_k$ has a degree-two
divisor over $\Q$.
\end{lemma}
\begin{proof}
Let $D\mid P_s$ with $\deg D=2$. Composing the isomorphisms of
Lemma~\ref{lem:algebras} with the quotient onto $\Q[x]/(D)$ gives a
surjection from $\Q[v]$ to a two-dimensional $\Q$-algebra, whose kernel
contains $g_k$. Its kernel is a principal ideal $(E)$, and
$\dim_{\Q}\Q[v]/(E)=\deg E=2$. Hence $E\mid g_k$ with $\deg E=2$.
This argument also applies when $D$ splits or has a repeated root.
\end{proof}

\section{The integral obstruction}
For positive odd $p,q$, put
\[
d=\frac{p^2+q^2}{2}\in\Z.
\]
Odd squares are $1$ modulo $4$, so $d$ is odd. Moreover
\[
k=\left(\frac{p^2-q^2}{p^2+q^2}\right)^2,
\qquad 1-k=\frac{p^2q^2}{d^2}.
\]
Multiplication by $d^2$ gives the integral quintic
\begin{equation}\label{eq:F}
F(v)=d^2g_k(v)
=d^2(v^5+v^4-2v^3-2v^2)+2p^2q^2.
\end{equation}

\begin{lemma}\label{lem:modular}
The polynomial $F$ in~\eqref{eq:F} has no degree-two divisor over $\Q$.
\end{lemma}
\begin{proof}
Suppose it has one. Gauss's lemma gives a factorization $F=A B$ in $\Z[v]$
with $\deg A=2$ and $\deg B=3$. There is no primitivity assumption on $F$:
clear denominators in a rational divisor, take its primitive part, and apply
Gauss's lemma to this primitive divisor. The complementary integral factor
absorbs any content of $F$.

The leading coefficients of $A$ and $B$ multiply to $d^2$, which is odd.
Their reductions modulo two therefore retain degrees two and three.
But in $\F_2[v]$,
\[
\overline F=v^4(v+1).
\]
Any divisor of $v^4(v+1)$ with nonzero constant coefficient is coprime to
$v^4$, so it divides $v+1$ and has degree at most one. Consequently both
$\overline A$ and $\overline B$ have zero constant coefficient. Thus $A(0)$
and $B(0)$ are even, forcing $4\mid F(0)$. This contradicts
\[
F(0)=2p^2q^2\equiv2\pmod4.\qedhere
\]
\end{proof}

\begin{proof}[Proof of Theorem~\ref{thm:main}]
Positivity and $p\ne q$ give $s>0$ and $s\ne1$.
A degree-two divisor of $P_s$ would give one of $g_k$ by
Lemma~\ref{lem:degree}, and hence one of $F=d^2g_k$.
Lemma~\ref{lem:modular} rules this out.
\end{proof}

\clearpage
\section{Formal verification and remaining work}
The final declaration is
\begin{center}\small\ttfamily
NumberTheoryResearch.BatchA.Cuboid.\\odd\_odd\_P\_no\_quadratic.
\end{center}
Its parameters are $p,q:\Z$, proofs that they are odd, positive, and
unequal, and $D:\Q[X]$ with natural degree two. Its conclusion is the
negation of $D\mid P_{(p/q)^2}$. The use of natural degree adds no zero
polynomial exception, because the zero polynomial has natural degree zero.

The sources use Lean 4.35.0 and mathlib commit
\begin{center}\small\ttfamily
e91b128082cc221d6432a6a5aadcd3243af8e76c.
\end{center}
The transitive axiom report for the theorem is exactly
\begin{center}\small\ttfamily
propext,\quad Classical.choice,\quad Quot.sound.
\end{center}
No admitted proof, custom axiom, or native evaluation axiom is used.
From the repository root, a fresh compilation and axiom audit runs with
\begin{quote}\small\ttfamily
scripts/with\_lean.sh lake env python3\\
research/parent/cuboid-formal/verify.py
\end{quote}
The command is a single shell line. The verifier compiles all eight source
modules into a new temporary directory, so it does not rely on previously
compiled research files. The public source repository is
\href{https://github.com/NoahWLono/GPT-6-Astra-Number-Theory}{GPT-6 Astra Number Theory}.

The parity argument is a classical local factor-degree obstruction.
Focused searches and inspection of the cited cuboid papers did not identify
this exact odd/odd restriction, but that does not establish originality.
The project therefore retains it as a verified research restriction with
novelty unverified, rather than counting it as a solved open problem.
The recent quadratic-factor paper~\cite{asiryan-rathbun} gives computational
evidence and a conditional conclusion for the full parameter range.
The present proof removes the odd/odd case from our remaining target;
it supplies no conclusion for the opposite-parity case.

\begin{thebibliography}{9}
\bibitem{sharipov}
R. Sharipov, \emph{Perfect cuboids and irreducible polynomials},
arXiv:1108.5348v1 (2011),
\url{https://arxiv.org/abs/1108.5348v1}.
\bibitem{sharipov-structural}
R. Sharipov, \emph{A note on the second cuboid conjecture. Part I},
arXiv:1201.1229v1 (2012),
\url{https://arxiv.org/abs/1201.1229v1}.
\bibitem{asiryan-rathbun}
V. Asiryan and R. L. Rathbun,
\emph{Computational Evidence Against Quadratic-Cubic Factorization for the
Second Cuboid Quintic}, arXiv:2601.07899v2 (2026),
\url{https://arxiv.org/abs/2601.07899v2}.
\end{thebibliography}
\end{document}
